\section{Introduction}
\label{sec:intro}

Potentiostats built around Analog Devices' AD5940/AD5941 electrochemical
analog front end (AFE) have become a recognizable instrument class within
the open-source, low-cost hardware literature: FreiStat~\cite{bill2023electrochemical},
HELPStat~\cite{bautista2025helpstat}, GaneStat~\cite{anshori2025ganestat},
and Vamos et al.'s HunStat2~\cite{vamos2025hunstat2} all use it as their
AFE, in the same lineage as discrete op-amp designs such as
CheapStat~\cite{rowe2011cheapstat}. A commercial AFE moves most of the
burden of getting a potentiostat's analog performance right -- transimpedance
gain accuracy, switch-matrix leakage, reference-voltage stability -- onto a
part that Analog Devices manufactures, tests, and characterizes at
volume~\cite{devices2024ad5940}. It does \emph{not} remove the burden of
driving that part correctly. The AD5940 exposes two independent,
purpose-built excitation/sensing paths -- a Low-Power (LP) loop for DC and
near-DC potentiostat operation, and a High-Speed (HS) loop for AC/impedance
excitation -- and firmware that configures the wrong one for a given
method, or that misinterprets the numeric convention of the raw code the
chip returns, produces code that compiles cleanly and runs without
crashing while still being physically wrong. Two design concerns recur for
an instrument built this way: (i) the firmware architecture needed to
correctly operate the AFE's two excitation/sensing paths for each
electrochemical method, and (ii) the software architecture needed to keep
that firmware maintainable as more methods and features are added. This
paper focuses its own diagnosis and fix on the first concern. A wrong
signal path or a wrong ADC convention
tends to produce output that is visibly, unmistakably broken -- flat,
noisy, or discontinuous -- which makes it easy to \emph{notice} but not
necessarily easy to \emph{root-cause} without a correctly-working
reference to compare against. This is a different failure mode from the
ones the two companion manuscripts on this instrument document: a
constant-gain calibration error is numerically well-behaved and passes
every linearity check~\cite{clement2026linearity}, and the software
architecture concern above is, on its own, a maintainability and
future-defect risk that need not affect any single
measurement~\cite{clement2026architecture} -- except, as
Section~\ref{sec:oop} shows, when the same duplication that makes an
architecture hard to maintain also means a single correctness bug exists
in multiple literal copies at once.

This paper reports two such defects, found in our firmware's
chronoamperometry (CA), square-wave voltammetry (SWV), and differential
pulse voltammetry (DPV) classes, by the most direct diagnostic method
available for a commercial-AFE design: comparing our implementation, file
by file and register by register, against the performance and
register-level behavior of Analog Devices' own AD5940 evaluation-kit
reference examples for the same class of
measurement~\cite{ad5940_examples_github} -- the same evaluation-kit code
this firmware, and the instrument as a whole
(Section~\ref{sec:instrument-design}), is built from. Prior internal testing on this
project had reported that CA and SWV produced ``flat, noisy scatter with no
discernible decay or peak'' against a resistive/capacitive dummy cell, while
cyclic voltammetry (CV) and open-circuit potential (OCP) -- which use a
different, independently-derived code path -- worked normally. That
asymmetry was itself a clue: whatever was wrong was specific to the three
affected classes' own signal-path configuration, not a property of the AFE
or the board.

\textbf{Contributions.} This paper makes the following contributions:
\begin{enumerate}
\item A structured audit of the firmware's hardware architecture, software
  architecture, and object-oriented design maturity
  (Section~\ref{sec:instrument-design}), using the actual source code as
  evidence: of eight C++ classes, only two exhibit non-trivial
  encapsulation, no inheritance or polymorphism is used anywhere in the
  codebase, and three method classes contain verbatim-duplicated
  measurement-configuration code.
\item A structured, file-by-file comparison of this project's CA/SWV/DPV
  firmware against Analog Devices' own \texttt{AD5940\_ChronoAmperometric},
  \texttt{AD5940\_SqrWaveVoltammetry}, and \texttt{AD5940\_Ramp} reference
  examples~\cite{ad5940_examples_github}, using the vendor's own code as
  ground truth rather than the presence of a \texttt{class} keyword or a
  plausible-looking register write.
\item Diagnosis and correction of two concrete firmware defects that
  together fully explain a previously-reported CA/SWV data-quality
  failure: (i) the three classes configured the AD5940's High-Speed
  loop where every vendor DC-method reference instead configures the
  Low-Power loop, and (ii) the ADC conversion routine used the wrong
  zero-point convention for the chip's offset-binary signed codes -- and
  the observation that, because of contribution 1's duplication, each had
  to be found and fixed three times over, once per copy.
\item An explicit classification of every other difference the comparison
  surfaced into \emph{deliberate} design choices appropriate to this
  instrument's live, host-driven use case, versus \emph{genuine,
  unreconciled} gaps -- including one gap this paper does not itself close
  (on-board RTIA live-calibration) but that the companion metrology
  manuscript's traceable protocol independently
  addresses~\cite{clement2026linearity}.
\item Hardware-validated performance evidence, using the same firmware
  baseline and traceable dummy-cell protocol validated in the two
  companion manuscripts~\cite{clement2026linearity,clement2026architecture},
  that the corrected signal chain produces correctly-signed,
  signal-dependent output for CA, CV, SWV, and DPV, and that its cyclic
  voltammetry is independently accurate ($R^2\geq0.999994$,
  $+0.428\,\%$ recovered-resistance error), not merely self-consistent.
\end{enumerate}

This paper's own audit (contribution 1) is deliberately scoped to what is
needed to motivate and contextualize its signal-path/ADC-convention fix:
it documents the duplication and the encapsulation gap as found, but does
not itself restructure the firmware. That structural refactor -- collapsing
the duplicated code into a shared base class, removing confirmed-dead
subsystems, and hardware-measuring that the refactor does not change any
method's output -- is reported in full, and separately validated, in the
companion architecture manuscript~\cite{clement2026architecture}; the
gain-calibration traceability protocol referenced throughout is reported
in full in the companion metrology manuscript~\cite{clement2026linearity}.
Sensor development, analyte concentration measurement, and
limit-of-detection studies remain explicitly out of scope for the whole
instrument family.
